\clearpage
\VerdierPhysicalPageAnchor{0177}%
\begingroup\small\setlength{\parskip}{0pt}\setlength{\emergencystretch}{1.5em}
\setlength{\abovedisplayskip}{4pt plus 2pt minus 2pt}%
\setlength{\belowdisplayskip}{4pt plus 2pt minus 2pt}%

\VerdierSemanticAnchor{ch3d-1-2-1-truncations-cont}
\noindent（相应地，
\[
\begin{aligned}
i^q&=\operatorname{id},&&q<n-1,\\
i^{n-1}&\colon\operatorname{Ker}(d_X^{n-1})\longrightarrow X^{n-1}
  &&\text{（典范态射）},\\
i^q&=0,&&q\geq n
\end{aligned}
\]
）。

注意，态射 $p$（相应地，$i$）在次数 $j\geq n$（相应地，$j<n$）的
上同调对象上诱导同构，并且当 $j<n$ 时 $H^j(X(n,+\infty))=0$
（相应地，当 $j\geq n$ 时 $H^j(X(-\infty,n))=0$）。利用这个构造，
命题 (1.2.1) 的所有断言都立即得到证明。详细证明留给读者。

\VerdierSemanticAnchor{ch3d-1-2-2}
\medskip\noindent\textbf{定义 1.2.2.}\quad
三角范畴
\[
D(\mathcal A)=K(\mathcal A)/\operatorname{Ac}(\mathcal A)
=K(\mathcal A)[\operatorname{Qis}(\mathcal A)^{-1}]
\]
称为 $\mathcal A$ 的\emph{导出范畴}，记为 $D(\mathcal A)$
（第二章，2.2.6 与 2.2.10）。

由定理 (2.2.6)，典范上同调函子
$H\colon K(\mathcal A)\longrightarrow\mathcal A$ 唯一地分解为一个仍用相同
记号表示的上同调函子：

\VerdierSemanticAnchor{ch3d-1-2-2-cohomology}
\begin{equation}
H\colon D(\mathcal A)\longrightarrow\mathcal A,
\tag{1.2.2.1}
\end{equation}
它仍称为\emph{典范上同调函子}。典范商函子记为
$Q\colon K(\mathcal A)\longrightarrow D(\mathcal A)$。

\VerdierSemanticAnchor{ch3d-1-2-3}
\medskip\noindent\textbf{定理 1.2.3.}\quad
\textit{a) 下图中的自然函子：}

\VerdierSemanticAnchor{ch3d-1-2-3-a-diagram}
\begin{center}
  \resizebox{0.80\textwidth}{!}{%
    $\begin{tikzcd}[ampersand replacement=\&,column sep=4.5em,row sep=3.0em]
      \& K^{,+}(\mathcal A)/\operatorname{Ac}(\mathcal A)
        \arrow[dr] \& \\
      K^{,b}(\mathcal A)/\operatorname{Ac}(\mathcal A)
        \arrow[ur] \arrow[dr]
        \& \& D(\mathcal A) \\
      \& K^{,-}(\mathcal A)/\operatorname{Ac}(\mathcal A)
        \arrow[ur] \&
    \end{tikzcd}$%
  }
\end{center}
\endgroup
